\ifx\danslelivre\undefined
\documentclass[10pt]{../fichiers-configuration-latex/smfart}
\usepackage{ifthen}
\newboolean{english}
\newboolean{francais}
\setboolean{english}{false}
\setboolean{francais}{true}
\usepackage{palatino,euler}
\input{../fichiers-configuration-latex/paquets}
\selectlanguage{francais}
\input{../fichiers-configuration-latex/lemme-proposition-theoreme}
\input{../fichiers-configuration-latex/numerotation}
\input{../fichiers-configuration-latex/macros}
\input{../fichiers-configuration-latex/ucs_manquants}
\input{../.cv}


\externaldocument{../04-Cohen-Gabber/Cohen-Gabber}
\externaldocument{../13-theoreme-de-finitude/finitude}

\begin{document}

\begin{center}
V. Algébrisation partielle\\
Fabrice Orgogozo
\end{center}
\version

\else
\setboolean{english}{false}
\setboolean{francais}{true}
\selectlanguage{francais}
\chapter{Algébrisation partielle}
\addtocontents{toc}{par Fabrice Orgogozo}
\begin{center}
Fabrice Orgogozo
\end{center}
\fi


\section{Préparatifs (rappels)}

\subsection{Le théorème de préparation de Weierstraß}
On trouvera dans \ac{vii}{3}{7-8}{} une démonstration du théorème
suivant.

\begin{théorème2}\label{5-Weierstrass}
Soient $A$ un anneau local séparé complet d'idéal maximal $𝔪$,
$d>0$ un entier et $f∈A[[\sous{X},T]]$ une série formelle,
où l'on pose $\sous{X}=(X₁,\dots,X_d)$.
\begin{enumerate}
\item Soit $ρ$ un entier naturel tel que
$f$ soit \emph{$ρ$-régulière} relativement à $T$, c'est-à-dire congrue à $(u∈A[[T]]^×)\cdot
T^ρ$ modulo $(𝔪,\sous{X})$. Alors, pour tout $g∈A[[\sous{X},T]]$, il existe un unique
couple $(q,r)∈A[[\sous{X},T]]×A[[\sous{X}]][T]$ tel que
$g=qf+r$ et $\deg_T(r)<ρ$. De plus, il existe
un unique polynôme $P=T^ρ+∑_{i<ρ} p_i T^i$, où
les coefficients $p_i$ appartiennent à $(𝔪,\sous{X})A[[\sous{X}]]$,
et une unité $u∈A[[\sous{X},T]]^×$ tels que $f=uP$.
\item Si $f$ est \emph{non nulle modulo $𝔪$},
il existe un entier naturel $ρ$ et un automorphisme
$A[[T]]$-linéaire $c$ de $A[[\sous{X},T]]$, tel que
$c(X_i)=X_i+T^{N_i}$ ($N_i>0$) et la série
entière $c(f)$ soit $ρ$-régulière.
\end{enumerate}
\end{théorème2}

\subsubsection{}Signalons que l'on peut satisfaire la condition (ii)
simultanément pour un nombre \emph{fini} d'éléments :
cf. \emph{loc. cit.}, n°7, lemme 2 où l'on considérera
un produit (fini) de séries formelles. 

Nous ferons usage de la propriété suivante des
polynômes comme en (ii) ci-dessus. 

\begin{lemme2}\label{5-lemme complétion}
Soient $B$ un anneau local complet nœthérien et $P∈B[X]$ un polynôme
de la forme $X^ρ+∑_{i<ρ}b_i X^i$, où $b_i∈\MM_{B}$ et
$ρ>0$. Alors, le complété $(P)$-adique de $B\{X\}$ s'identifie à
$B[[X]]$.
\end{lemme2}

Rappelons que $B\{X\}$ désigne l'hensélisé en l'origine de l'anneau $B[X]$.
Un polynôme $P$ comme ci-dessus est parfois dit \emph{de Weierstraß}.

\begin{démo}
Soient $N$ un entier naturel et $Q=P^N$.
Il résulte de \ref{5-Weierstrass} (i), que l'anneau quotient
$B[[X]]/(Q)$ est isomorphe comme $B$-module à $B[X]/(X^{\deg(Q)})$
et en particulier fini sur $B$. 
Par fidèle platitude du morphisme $B\{X\}→B[[X]]$, on a
$QB[[X]]∩B\{X\}=QB\{X\}$ de sorte que le $B$-morphisme
$B\{X\}/Q→B[[X]]/Q$ est \emph{injectif}. L'anneau $B\{X\}/Q$
est donc également fini sur l'anneau complet $B$ ;
il est donc isomorphe à son complété $B[[X]]/Q$. En faisant tendre
$N$ vers l'infini, on en déduit que le séparé-complété $(P)$-adique de $B\{X\}$ est 
donc isomorphe à celui de $B[[X]]$ ; ce dernier est isomorphe à $B[[X]]$ 
puisque $\deg(P)>0$.
\end{démo}

\subsection{Le théorème d'algébrisation d'Elkik}

\begin{définition2}
Une paire $(\Spec(C),V(J))$ est dite \emph{hensélienne}
si pour tout polynôme $f∈C[T]$, toute racine $β$ de $f$
dans $C/J$ telle que $f'(β)$ soit une unité de $C/J$
se relève en une racine dans $C$.
\end{définition2}

\subsubsection{}Il résulte de \cite{Caracterisation@Crepeaux}, prop. 2, (iii)⇒(i)
et \cite{Anneaux@Raynaud}, chap. XI, §2, prop. 1, 4)⇒1)
(voir aussi \cite{K-theory@Gabber}, déf. p. 59) que la
définition précédente est équivalente à la définition
\ega{iv}{18.5.5} (relèvement d'ouverts-fermés).
Remarquons que la notion de \emph{racine simple} introduite
dans la définition est plus forte que celle de \a{iv}{2}{1}{, déf. 1} et que le relèvement
ci-dessus est nécessairement unique. Notons également
que l'idéal $I$ est nécessairement contenu
dans le radical de Jacobson de l'anneau $C$.

\begin{lemme2}\label{5-paire hensélienne}
Soient $C$ un anneau \emph{local} hensélien d'idéal maximal
$𝔪$, et $J⊆𝔪$ un idéal. La paire $(\Spec(C),V(J))$ est hensélienne.
\end{lemme2}

En particulier, pour $B$ et $P$ comme dans le lemme \ref{5-lemme
complétion}, la paire $(\Spec(B\{X\}),V(P))$  est hensélienne.

\begin{démo}
Soient $f$ et $β$ comme ci-dessus. L'anneau $C$ étant local hensélien,
l'image $γ$ de $β$ dans le corps résiduel $C/𝔪$ se relève en
une racine $α$ de $P$. Notons $β'$ son image dans $C/J$ et vérifions que $β=β'$.
Remarquons tout d'abord que puisque $P'(α)$ est une unité de $C$,
$P'(β')$ est une unité de $C/J$. De plus, l'égalité
$P(β)=P(β')+(β-β')P'(β')+(β-β')²b$ où $b∈B/J$ se réduit
à $β-β'=(β-β')²\frac{-b}{P'(β')}$ ; si l'on pose $x=β-β'$, on a donc $x(1-ax)=0$
pour un $a∈C/J$. Comme $x$ appartient à $𝔪$ (car $β$ et $β'$ ont
pour image $γ$ dans $C/𝔪$), on a $x=0$.
\end{démo}

Terminons ces rappels par l'énoncé du théorème d'algébrisation de Renée Elkik (\cite{Solutions@Elkik}, théorème 5).

\begin{théorème2}\label{5-Elkik}
Soient $(X=\Spec(A),Y=V(I))$ une paire hensélienne avec $A$ nœthérien,
et $U$ le sous-schéma ouvert complémentaire de $Y$ dans $X$.
Notons $X_{\chap{Y}}$ le complété de $X$ le long de $Y$, $\chap{Y}$
le fermé correspondant à $Y$ et $\chap{U}$ son complémentaire dans 
$X_{\chap{Y}}$.
Le foncteur $X'↦X'×_X X_{\chap{Y}}$ induit une
équivalence de catégories entre la catégorie des
$X$-schémas finis, étales sur $U$, et la catégorie
des $X_{\chap{Y}}$-schémas finis, étales sur $\chap{U}$. 
\end{théorème2}

\section{Algébrisation partielle en égale caractéristique}\label{5-égale caractéristique}

\subsection{Énoncé}

\subsubsection{}\label{5-déf algébrisation partielle}
Soient $A$ un anneau local nœthérien complet et $\{I_e\}_{e∈E}$ une
collection d'idéaux de $A$. On dit que la paire $(A,\{I_e\}_{e∈E})$ est \emph{partiellement algébrisable} 
s'il existe un anneau local nœthérien complet $B$ de
dimension strictement inférieure à celle de $A$, une
$B$-algèbre de type fini $C$, un idéal maximal $𝔫$ au-dessus
de l'idéal maximal de $B$, et un isomorphisme
$A\isononcan \chap{C_{𝔫}}$ tel que les idéaux $I_e$ ($e∈E$) proviennent
d'idéaux de $C_{\got{n}}$.

\begin{théorème2}\label{5-algébrisation partielle égale caractéristique}
Soit $A$ un anneau local nœthérien complet réduit d'égale caractéristique
qui ne soit pas un corps. Alors, $A$ muni d'un ensemble \emph{fini} quelconque
d'idéaux est partiellement algébrisable.
\end{théorème2}

\subsection{Démonstration}
\subsubsection{}Soient $X=\Spec(A)$ et $I_1,\dots,I_n⊆A$ comme dans
l'énoncé. Il résulte de la définition \ref{5-déf algébrisation partielle}
que si un idéal $I$ de $A$ est de la forme $J₁∩\cdots∩J_r$ 
et que les $J_i$ sont partiellement algébrisables, 
l'idéal $I$ l'est également. D'après le théorème de
décomposition primaires des idéaux, on peut supposer les $I_i$
\emph{primaires}.

\subsubsection{}Notons $k$ le corps résiduel de $A$ et $d>0$ sa dimension. 
D'après \refext{IV}{4-Cohen-Gabber égale} si $X$ est équidimensionnel
ou bien \refext{IV}{4-Cohen-Gabber équivariant} (avec $G=\{1\}$) dans
le cas général, il existe un morphisme fini génériquement
étale $π:X→X₀$, où $X₀=\Spec(k[[t₁,\dots,t_d]])$.

\subsubsection{}Soit $I$ l'un des $I_i$. Deux cas se présentent.
\begin{enumerate}
\item $\dim(A/I)=d$. L'idéal $I$ est donc un idéal premier
minimal de $A$. 
\item $\dim(A/I)<d$. L'image de $V(I)$ dans $X₀$
est donc de dimension au plus $d-1$ donc contenue dans un fermé
$V(g_I)$ où $g_I∈A₀=Γ(X₀,𝒪_{X₀})-\{0\}$.
\end{enumerate}

Soient $g=∏_{I_i} g_{I_i}$ où $I_i∈\{I_1,\dots,I_n\}$ 
parcourt le sous-ensemble des idéaux du second type,
et $f∈A₀-\{0\}$ telle que le lieu de ramification de
$π$ soit contenu dans $V(f)$. Posons $h=gf$.
D'après \ref{5-Weierstrass} (ii) et (i),
quitte à changer de base par un automorphisme (c'est-à-dire changer les
coordonnées), on peut supposer que $h$ est un polynôme
unitaire en $t_d$. Considérons le sous-anneau $\gtilde{A₀}=k[[t₁,\dots,t_{d-1}]]\{t_d\}$
de $A₀$. Il est hensélien et contient $h$. 
D'après les lemmes \ref{5-lemme complétion} et \ref{5-paire hensélienne}, 
la paire $(\gtilde{X₀}=\Spec(\gtilde{A₀}),V(h))$ est hensélienne. On est donc en
mesure d'appliquer le théorème \ref{5-Elkik} et d'en déduire qu'il
existe un diagramme cartésien :
$$
\xymatrix{X \ar[d] \ar[r] & \gtilde{X} \ar[d] \\
X₀ \ar[r] & \gtilde{X₀}
}
$$
où la flèche verticale de gauche est, par hypothèse, étale hors de 
$V(h)$ et les flèches horizontales sont des morphismes de complétion
(à la fois pour la topologie $h$-adique et celle définie par leurs
idéaux maximaux respectifs).

Les idéaux $I$ du premier type (c'est-à-dire premier minimaux)
se descendent à $\gtilde{X}$ d'après le lemme suivant.

\begin{lemme2}\label{5-premier minimal}
Soit $B$ un anneau local hensélien quasi-excellent de complété
noté $\chap{B}$. Tout idéal premier minimal de $\chap{B}$ provient par
image inverse d'un idéal premier minimal de $B$.  
\end{lemme2}

\begin{démo}
Par restriction à l'adhérence du point générique du fermé,
il suffit de démontrer que le complété d'un intègre hensélien
quasi-excellent est intègre. Ce fait est bien connu et résulte
d'ailleurs immédiatement du théorème d'approximation de
Popescu, appliqué à l'équation \mbox{$xy=0$}.
\end{démo}

\subsubsection{}Quant aux idéaux $I$ du second type, il suffit d'observer que 
chaque $V(I)$ est fini sur $\Spec(k[[t₁,\dots,t_{d-1}]])$,
donc sur $\gtilde{X}$, et d'appliquer 
le

\begin{lemme2}\label{5-fini par complétion}
Soient $B$ un anneau local nœthérien, $J⊆𝔪_B$ un idéal, et 
$\chap{B}$ le complété $J$-adique de $B$.
Tout quotient de $\chap{B}$ fini sur $B$ se descend à $B$.
\end{lemme2}

\subsubsection{}Admettons momentanément ce lemme et achevons la démonstration 
de \ref{5-algébrisation partielle égale caractéristique}. Comme on l'a vu,
les idéaux $I_1,\dots,I_n$ proviennent d'idéaux de $\gtilde{A}$.
Cet anneau est fini — donc \emph{a fortiori} de présentation finie
par nœthérianité — sur l'anneau $k[[t₁,\dots,t_{d-1}]]\{t_{d}\}$. Ce dernier est
l'hensélisé en l'origine de $k[[t₁,\dots,t_{d-1}]][t_d]$ ;
il est isomorphe à la colimite filtrante d'anneaux de type fini sur
$k[[t₁,\dots,t_{d-1}]][t_d]$ donc sur l'anneau $B=k[[t₁,\dots,t_{d-1}]]$.
La conclusion résulte de \ega{iv}{8.8.2} qui assure l'existence d'une
$B$-algèbre $C$ comme en \ref{5-déf algébrisation partielle}
dont proviennent les idéaux $I_i$.

\subsubsection{}Revenons à la démonstration du lemme \ref{5-fini par
complétion} ci-dessus.
Soit $I⊆\chap{B}$ tel que $\chap{B}/I$ soit fini sur $B$.
Quitte à remplacer $B$ par $B/\Ker(B→\chap{B}/I)$,
c'est-à-dire $\Spec(B)$ par l'image schématique de $V(I)$, on peut supposer
$B→\chap{B}/I$ injectif, c'est-à-dire $V(I)→\Spec(B)$ schématiquement
dominant. Le $B$-module $\chap{B}/I$ étant fini, la topologie
$J$-adique sur $\chap{B}/I$ induit la topologie $J$-adique sur $B$.
Puisque l'application $B→\chap{B}/I$ est injective, d'image dense,
et continue, il en résulte que $\chap{B}/I$ est le
\emph{séparé-complété} de $B$ pour la topologie $J$-adique. On a donc
$I=(0)$ ; il se descend tautologiquement à $B$.

\section{Algébrisation partielle première à $ℓ$ en caractéristique mixte}

\subsection{Énoncé}

\begin{théorème2}
\label{5-algébrisation partielle mixte raffinée}
Soient $A$ un anneau local nœthérien complet normal de caractéristique mixte
$(0,p)$ de dimension $d≥2$ et $ℓ≠p$ un nombre premier. Il existe
un morphisme injectif fini $A→A'$ de degré générique premier à $ℓ$,
où $A'$ est un anneau normal intègre dont toute famille \emph{finie}
d'idéaux est partiellement algébrisable.
\end{théorème2}

\begin{remarque2}Signalons qu'il suffit pour démontrer \refext{XIII}{13-constructibilité}
d'établir la variante affaiblie de l'énoncé précédent selon laquelle
— en reprenant les notations de \ref{5-déf algébrisation partielle} —
tout fermé rare de $\Spec(A')≃\Spec(\chap{C_{𝔫}})$
est ensemblistement contenu dans l'image inverse d'un \emph{diviseur} de
$\Spec(C_𝔫)$.
\end{remarque2}

\subsubsection{Reformulation géométrique}
\label{5-algébrisation partielle mixte raffinée géométrique}
Soit $X$ un schéma local nœthérien complet de caractéristique mixte,
de dimension $d ≥ 2$ et soit $\{Z_i\}_{i ∈ I}$ une famille finie de
fermés de $X$. Pour tout nombre premier $ℓ$ inversible sur $X$,
il existe un diagramme
\[
\xymatrix{X & X ′ \ar[r]^{a} \ar[l]_{π} & Y \ar[d]^f  \\
& & S }
\]
où :
\begin{itemize}
\item $S$ est un schéma nœthérien régulier complet de caractéristique
mixte et de dimension $d-1$ ;
\item $X ′$ est un schéma local normal ;
\item $π$ est un morphisme fini de degré générique premier à $ℓ$ ;
\item $f$ est un morphisme de type fini ;
\item $a$ induit un isomorphisme entre $X ′$ et le complété de $Y$ en
un point fermé de la fibre spéciale de $f$,
\end{itemize}
et des fermés $F_i$ de $Y$ tels que $π^{-1}(Z_i)=a^{-1}(F_i)$ pour
tout $i ∈ I$.

\begin{remarque2}
Il découle de \ref{5-algébrisation partielle égale caractéristique}
que le résultat précédent est également vrai en égale caractéristique,
et que l'on peut alors supposer $X=X′$.
\end{remarque2}

\subsection{Démonstration}

\subsubsection{}Soient $X=\Spec(A)$ de dimension $d≥2$ et $ℓ$ comme
dans l'énoncé. D'après le théorème \refext{IV}{4-Cohen-Gabber-Epp équivariant},
il existe un diagramme commutatif
$$
\xymatrix{
V[[t_1,\dots,t_{d-1}]]=B₀ \ar[r] & B  & A'=\Fix_H(B) \ar[l]\\
& A \ar[u] \ar[ur] &
}
$$
où $V$ est le spectre d'un anneau de valuation discrète complet
d'idéal maximal $𝔪_V$, et $H$ est un $ℓ$-groupe agissant sur l'anneau
normal $B$, son sous-anneau $V$, et trivialement sur les variables $t_i$
($1≤i≤d-1$). De plus, $A→A'$ est une injection finie de rang générique premier à $ℓ$
et $π:\Spec(B)→\Spec(B₀)$ est fini, $p$-génériquement étale.

\subsubsection{}Nous allons montrer que toute famille finie d'idéaux de $A'$ est partiellement
algébrisable. Soit $I'₁,\cdots,I'_n$ une telle famille, que l'on peut 
supposer constituée d'idéaux primaires (cf. §\ref{5-égale caractéristique}).

\subsubsection{}Soit $I'$ l'un des $I_i$. Deux cas se présentent. 
\begin{enumerate}
\item $\dim(A'/I+p)=d-1$. Supposons $I'≠(0)$
et notons $𝔭$ l'idéal premier, nécessairement de hauteur un,
pour lequel $I'$ est primaire. Par hypothèse, l'idéal premier $𝔭$ contient $p$ ;
c'est un idéal premier minimal de $A'/p$. D'autre part, 
$A'$ est normal car $B$ l'est. Il résulte par exemple
de \cite{Algebre@Serre}, chap. \textsc{iii}, C, §1 que $I'$
est une \emph{puissance symbolique} de $𝔭$, c'est-à-dire
l'image inverse dans $A'$ d'une puissance de l'idéal (principal) $𝔭A'_𝔭$.
\item $\dim(A'/(I+p))<d-1$. L'image de $V(I)$ dans
$\Spec(V^H[[t₁,\dots,t_{d-1}]])$ est donc contenu dans un fermé
$V(g_{I'})$ où $g_{I'}∈V^H[[t₁,\dots,t_{d-1}]]$ est \emph{non nulle
modulo $𝔪_{V^H}$}.
\end{enumerate}

Soient $g=∏_{I'_i} g_{I'_i}$ où $I'_i∈\{I'_1,\dots,I'_n\}$ 
parcourt le sous-ensemble des idéaux du second type,
et $f∈V[[t₁,\dots,t_{d-1}]]-𝔪_{V}$ telle que le lieu de ramification de
$π$ soit contenu dans $V(f)$. Posons $h=gf$.
D'après le théorème de préparation (\ref{5-Weierstrass}) et l'observation lui faisant
immédiatement suite, on peut supposer que $h$ est un polynôme
de Weierstraß en $t_{d-1}$, \emph{$H$-équivariant}. (Rappelons que $H$ agit
trivialement sur les variables).
Comme en §\ref{5-égale caractéristique}, le morphisme $B₀→B$ se descend donc
d'après \ref{5-Elkik} en un morphisme
$\gtilde{B₀}→\gtilde{B}$, où $\gtilde{B₀}=V[[t_1,\dots,t_{d-2}]]\{t_{d-1}\}$. 
Le groupe $H$ préservant l'ouvert $D(f)$ de $\Spec(B₀)$, son action se
descend. Le diagramme ci-dessus se complète donc en un diagramme
$$
\xymatrix{
V[[t_1,\dots,t_{d-2}]]\{t_{d-1}\}=\gtilde{B₀} \ar[d] \ar[r] &
\gtilde{B} \ar[d] & \gtilde{A'}=\gtilde{B}^H \ar[l] \ar[d] \\
B₀ \ar[r] & B  & A' \ar[l]
}
$$
où les flèches verticales sont les morphismes de complétion
et les flèches horizontales sont finies.

\subsubsection{}Les idéaux $I'$ du second type se descendent de $A'$ à $\gtilde{A'}$
car $A'/I'$ est fini sur $V^H[[t_1,\dots,t_{d-2}]]$ donc \emph{a
fortiori} sur $\gtilde{A'}$ (cf. \ref{5-fini par complétion}).
Quant aux idéaux du premier type (puissances symboliques), il suffit d'appliquer
le lemme \ref{5-premier minimal} à la paire constituée de $\gtilde{A'}/p$ et de
son complété $A'/p$.
Comme en §\ref{5-égale caractéristique}, on utilise le fait que
$\gtilde{A'}$ soit fini — de type fini suffirait — sur $V^H[[t_1,\dots,t_{d-2}]]\{t_{d-2}\}$ pour
descendre, par passage à la limite, les idéaux à un anneau
de type fini sur $V^H[[t_1,\dots,t_{d-2}]]$.

\ifx\danslelivre\undefined
\bibliography{../fichiers-configuration-latex/bibliographie}
\bibliographystyle{../fichiers-configuration-latex/nomcomplet-date}
\end{document}
\fi
